\documentclass[../../../include/open-logic-section]{subfiles}

\begin{document}

\olfileid[pt]{sth}{spine}{zf}
\olsection{$\Z$ e $\ZF$: um marco}

Com a Fundação, atingimos outro marco importante. Consideramos as teorias $\Zminus$ e $\ZFminus$, que dissemos ser certas teorias «menos» alguma coisa. Essa coisa é a Fundação. Logo:

\begin{defn}
A teoria $\Z$ acrescenta a Fundação a $\Zminus$. Os seus axiomas são
Extensionalidade, União, Pares, Conjuntos das Partes, Infinito, Fundação, e todas
as instâncias do esquema de Separação.

A teoria $\ZF$ acrescenta a Fundação a $\ZFminus$. Por outras palavras, $\ZF$
acrescenta todas as instâncias de Substituição a~$\Z$.
\end{defn}

Ainda assim, uma pergunta pode ter surgido. Se a Regularidade é
equivalente sobre $\ZFminus$ à Fundação, e a justificação da Regularidade é clara, por que dar voltas e tomar
a Fundação como axioma básico, em vez da Regularidade? 

Deixando de lado razões históricas (quem formulou o quê e
quando), a razão essencial é que a Fundação pode ser apresentada sem
recorrer à definição dos~$V_\alpha$. Essa definição dependeu
de todo o trabalho na \olref[recursion]{sec}: foi preciso
provar a recursão transfinita, para mostrar que era legítima.
Mas a nossa prova da recursão transfinita usou o \emph{Esquema de Substituição}.
Logo, embora a Fundação e a Regularidade sejam equivalentes módulo $\ZFminus$,
não são equivalentes módulo $\Zminus$. 

De fato, a questão é mais grave do que esta observação simples sugere.
Embora ultrapasse o âmbito deste livro, verifica-se que tanto
$\Zminus$ como $\Z$ são demasiado fracas para definir os $V_\alpha$. Logo, se
se trabalha apenas em $\Z$, então a Regularidade (tal como a formulamos)
nem sequer faz \emph{sentido}. É por isso que o nosso axioma oficial é a
Fundação, e não a Regularidade. 

Daqui em diante, trabalharemos em $\ZF$ (salvo indicação em contrário), sem
mais comentários. 

\end{document}
